\section{Inpainting}
\label{sec:inpainting}

\iwh{This section should contain the theoretical motivation behind inpainting and the details of how it is applied \emph{in general}. I'm trying to work on this at the moment.}

\gscd{Need to be somewhat careful here, as it doesnt particularly follow on from the zerolatency filter method.
Motivate its use first, then go into details}

We have showed in the previous section, and in our previous work~\cite{CabournDavies:2024hea}, that the zero-latency filter method, used successfully for many years in ground-based low-latency analysis~\cite{Tsukada:2017cuf} is directly applicable to
solve the \ac{MBHB} premerger problem in LISA.

However, when considering the case of LISA, there are some issues with this method that warrant considering alternatives. The premerger analysis that we carried out with the zero-latency filter here analyzes 30 days of data to identify mergers up to 14 days before merger. LISA data is expected to have gaps \iwh{say a bit about expected cadence of gaps and cite some LISA design/white paper}. We do not have a method for dealing with these gaps using the zero-latency filter and therefore if there is a gap in the recorded LISA data in the days before a \ac{MBHB} merger it might result in a delay in identifying that the merger is about to happen.

Here we consider an alternative technique, which is also adapted from ground-based analyses. In particular the observation of the first binary neutron star merger, GW170817, was significantly hampered by the presence of a loud non-Gaussian artefact in the seconds before the merger. At the time of observation the glitch was manually excised from the data, in a process known as "gating". However, this gating might result in spectral leakage in the whitened strain, as described in~\cite{Zackay:2019kkv}, which can cause issues for detection pipelines and bias inference algorithms.

The method of ``inpainting'' for proposed in~\cite{Zackay:2019kkv} as an
alternative for safely removing non-Gaussianities or, equivalently, for analysing over data gaps. The full derivation is given in~\cite{Zackay:2019kkv,Capano:2021etf} but the basic idea is that we want the ``overwhitened'' strain to be 0s for the duration of the gap. If we consider the matched-filter between signal, $s$ and data $h$
%
\begin{equation}
    (h|s) = 4\mathrm{Re}\int_{f_\mathrm{low}}^{f_\mathrm{high}}\frac{\tilde{h}(f)\tilde{s}^*(f)}{S_n(f)} \mathrm{d} f,
\end{equation}
%
we can rewrite this as
%
\begin{equation}
    (h|s) = 4\mathrm{Re}\int_{f_\mathrm{low}}^{f_\mathrm{high}} \tilde{h}(f) \frac{\tilde{s}^*(f)}{S_n(f)} \mathrm{d} f,
\end{equation}
%
where $\tilde{s}^*(f) / S_n(f)$ is the ``overwhitened'' strain. The matched-filter is then the cross correlation between the overwhitened strain and the template waveform and therefore any 0s in the overwhitened strain will not contribute to the overall matched-filter. Choosing values that the data should take in the gap of the [unwhitened] strain such that the overwhitened strain is 0 boils down to a linear algebra problem.


Gaps in the data will cause significant issue in the \ac{LISA} data analysis\gscd{CITE}, and there is not a way to
implement the zero-latency filter method to account for multiple gaps in the data.
As a result, we consider the use of \textit{inpainting} - a data analysis technique used to fill in missing
segments of detector data~\cite{Zackay:2019kkv}.
\gscd{This is probably where we can really big up inpainting - the zerolatency filter method can't deal with data gaps.}

To describe inpainting, we consider matched filtering, and how the removal of data in each term
must be considered.

The basic \ac{SNR} calculation is
\begin{equation}
    \rho^2(t) = \frac{[(h|s)]^2}{(h|h)}
\end{equation}

where the inner product is
\begin{equation}
    (a|b)(t) = 4\mathrm{Re}\int_{f_\mathrm{low}}^{f_\mathrm{high}}\frac{\tilde{a}(f)\tilde{b}^*(f)}{S_n(f)} \mathrm{e}^{2i \pi ft} \mathrm{d} f,
\end{equation}

$\tilde{a}$ defines the Fourier transform, $\tilde{b}^*(f)$ is the complex conjugate 
of the Fourier transform, and $S_n(f)$ is the one-sided \ac{PSD} of the detector noise.

\subsection{$(h|s)$}
\label{subsubsec:h_s}
Dividing the data $\tilde{s}(f)$ by the square root of the \ac{PSD} - the \ac{ASD}, $\sqrt{S_n(f)}$ -
of the noise is often referred to as \textit{whitening} the data, as the spectral density of the data will
become flat by this procedure. By performing this step twice on the data, we produce \textit{overwhitened} data
\footnote{Also sometimes refered to as \textit{blued} data, particularly in \cite{Zackay:2019kkv}}.
With overwhitened data, we are able to use un-whitened templates in the inner product calculation.

The method of inpainting is to add values to the gap so that after the overwhitening procedure, data outside
the gap time is unaffected, but data within the gap becomes zero. If the data within the gap can be identically
zero after overwhitening, then the data from this time will not contribute to the \ac{SNR} calculation.

The challenge to calculate the values to be added to the gap is one of linear algebra, fully described
in~\cite{Zackay:2019kkv}, and implementated in \texttt{PyCBC} for~\cite{Capano:2021etf,Wang:2023ljx}, 

\subsection{$(h|h)$}
\label{subsubsec:h_h}


\subsection{Use as a Premerger Search}
\label{subsec:inpainting_premerger}
Inpainting does not require data either side in order to `bridge' a gap, and so
we can consider the missing data in the future to simply be a one-sided gap.

For ringdown analyses, once can use inpainting to remove data from the inspiral~\cite{Capano:2021etf,Wang:2023ljx},
so that the masses and spins can be inferred independently from the ringdown and inspiral, and cross-checked.
Here, we propose a similar strategy, but removing data from the opposite end; by inpainting and setting all
future data to zero, we can - similarly to~\cite{CabournDavies:2024hea} - implement a pre-merger search without
the requirement to wait for future data.
This solution works for LISA but not ground-based premerger searches \gscd{CITE!} as the speed requirements of
ground-based pre-merger analyses mean that we cannot wait for teh inpainting procedure in real time \gscd{word this better}